\documentclass[11pt]{article}

\usepackage[utf8]{inputenc}
\usepackage[T1]{fontenc}
\usepackage{lmodern}
\usepackage{geometry}
\usepackage{amsmath,amssymb,amsfonts,amsthm,mathtools}
\usepackage{bm}
\usepackage{enumitem}
\usepackage{microtype}
\usepackage{hyperref}
\usepackage{titlesec}
\usepackage{booktabs}
\usepackage{array}
\usepackage{setspace}
\usepackage{mathrsfs}
\usepackage{physics}

\geometry{margin=1in}
\hypersetup{colorlinks=true, linkcolor=black, urlcolor=blue, citecolor=black}
\setlist[enumerate]{topsep=0.4em,itemsep=0.35em}
\setlist[itemize]{topsep=0.3em,itemsep=0.25em}
\titleformat{\section}{\large\bfseries}{\thesection.}{0.5em}{}
\titleformat{\subsection}{\normalsize\bfseries}{\thesubsection.}{0.5em}{}
\setstretch{1.06}

\newcommand{\Aether}{\AE{}ther}
\newcommand{\Aflow}{\AE{}ther-flow}
\newcommand{\TheoryName}{\Aflow{} Interpretation of Relativity}
\newcommand{\TheoryProgram}{\Aether{} / \Aflow{} framework}

\newcommand{\CompletedMetric}{g^{\sharp}}
\newcommand{\MatterSpace}{\Psi}

\newcommand{\Car}{\mathrm{car}}
\newcommand{\Cur}{\mathrm{cur}}
\newcommand{\Hom}{\mathrm{hom}}

\newcommand{\ForSpace}[1]{\mathcal I_{#1}^{\mathrm{for}}}
\newcommand{\PrimShell}{\mathcal S_{\mathrm{prim}}}
\newcommand{\MetricOut}{\mathscr G_{\mathrm{out}}}
\newcommand{\MatterOut}{\mathscr T_{\mathrm{out}}}

\newcommand{\SubSectionPackage}{\mathfrak V^{\mathrm{subsec}}}
\newcommand{\MicroSectionPackage}{\mathfrak W^{\mathrm{microsec}}}
\newcommand{\MicroSectionSpace}[1]{\mathcal N_{#1}^{\mathrm{microsec}}}
\newcommand{\MicroSectionBody}[1]{\mathcal B_{#1}^{\mathrm{microsec}}}
\newcommand{\MicroSectionFunctional}[1]{\mathscr E_{#1}^{\mathrm{microsec}}}
\newcommand{\MicroSectionShell}[1]{\mathcal Z_{#1}^{\mathrm{microsec}}}
\newcommand{\SubPreSectionCore}{\mathfrak U^{\mathrm{sec}}}

\newcommand{\NanoSectionPackage}{\mathfrak X^{\mathrm{nanosec}}}
\newcommand{\NanoSectionSpace}[1]{\mathcal P_{#1}^{\mathrm{nanosec}}}
\newcommand{\NanoSectionBody}[1]{\mathcal B_{#1}^{\mathrm{nanosec}}}
\newcommand{\NanoSectionProj}[1]{\Pi_{#1}^{\mathrm{nanosec}}}
\newcommand{\NanoSectionFiber}[2]{\mathcal F_{#1}^{\mathrm{nanosec}}(#2)}
\newcommand{\NanoSectionMicroMin}[1]{\mathfrak n_{#1}^{\mathrm{nanosec}}}
\newcommand{\NanoSectionFunctional}[1]{\mathscr N_{#1}^{\mathrm{nanosec}}}
\newcommand{\NanoSectionShell}[1]{\mathcal Z_{#1}^{\mathrm{nanosec}}}

\newcommand{\PicoSectionPackage}{\mathfrak Y^{\mathrm{picosec}}}
\newcommand{\PicoSectionSpace}[1]{\mathcal Q_{#1}^{\mathrm{picosec}}}
\newcommand{\PicoSectionBody}[1]{\mathcal B_{#1}^{\mathrm{picosec}}}
\newcommand{\PicoSectionProj}[1]{\Pi_{#1}^{\mathrm{picosec}}}
\newcommand{\PicoSectionFiber}[2]{\mathcal F_{#1}^{\mathrm{picosec}}(#2)}
\newcommand{\PicoSectionNanoMin}[1]{\mathfrak p_{#1}^{\mathrm{picosec}}}
\newcommand{\PicoSectionFunctional}[1]{\mathscr P_{#1}^{\mathrm{picosec}}}
\newcommand{\PicoSectionShell}[1]{\mathcal Z_{#1}^{\mathrm{picosec}}}

\newcommand{\PicoMinSectionCore}{\mathfrak Y^{\mathrm{pic,sec}}}

\newcommand{\FemtoSectionPackage}{\mathfrak Z^{\mathrm{femtosec}}}
\newcommand{\FemtoSectionSpace}[1]{\mathcal R_{#1}^{\mathrm{femtosec}}}
\newcommand{\FemtoSectionBody}[1]{\mathcal B_{#1}^{\mathrm{femtosec}}}
\newcommand{\FemtoSectionProj}[1]{\Pi_{#1}^{\mathrm{femtosec}}}
\newcommand{\FemtoSectionFiber}[2]{\mathcal F_{#1}^{\mathrm{femtosec}}(#2)}
\newcommand{\FemtoSectionPicoMin}[1]{\mathfrak f_{#1}^{\mathrm{femtosec}}}
\newcommand{\FemtoSectionFunctional}[1]{\mathscr F_{#1}^{\mathrm{femtosec}}}
\newcommand{\FemtoSectionShell}[1]{\mathcal Z_{#1}^{\mathrm{femtosec}}}

\newtheorem{definition}{Definition}
\newtheorem{proposition}{Proposition}
\newtheorem{remark}{Remark}
\newtheorem{corollary}{Corollary}
\newtheorem{theorem}{Theorem}

\title{\textbf{The \TheoryName{}}\\[0.4em]
\large Primitive-Reservoir Observer-Reduction\\
\large Nano-Section-Generating Substrate Pico-Section Dynamics\\
\large from Pico-Section-Generating Substrate Femto-Section Dynamics}
\author{}
\date{}

\begin{document}

\maketitle

\begin{abstract}
The current active-primary theorem already moved the live burden below the observer-relevant pico-section minimizer-section core itself. That theorem derived the current minimizer-section core canonically from the deeper current nano-section-generating substrate pico-section dynamics package \(\PicoSectionPackage\). The accompanying renewed routing note then fixed the next honest move: either derive or justify \(\PicoSectionPackage\) itself from still more primitive substrate structure, or prove a genuinely smaller theorem-level collapse strictly inside that deeper package.

The present note takes the direct derivation route. For each sector it introduces \emph{pico-section-generating substrate femto-section dynamics}: a compact convex femto-section body over the recorded pico-section body, an affine projection to that fixed body, a nonnegative smooth femto-section functional with unique fiberwise minimizers, and an exact femto-section shell projecting diffeomorphically to the recorded exact pico-section shell. The current pico-section functional is then no longer an unexplained primitive datum. It is the effective fiberwise minimum value induced by that deeper femto-section law, and the current exact pico-section shell is the projected exact shell of the deeper package.

The benchmark boundary remains fixed throughout. The one operative metric is still exactly \(\CompletedMetric\), the ordinary matter route is unchanged, there is no second operative metric, the low-energy matter law is not enlarged, and no stronger promotion language is licensed. The positive result is therefore narrow but benchmark facing: the current nano-section-generating substrate pico-section dynamics are no longer left as the primitive stopping point on the active line. The remaining honest burden moves one level deeper again, to the pico-section-generating substrate femto-section dynamics themselves or to a sharper collapse strictly inside that deeper package.
\end{abstract}

\section{Introduction}

The active derivational program of \TheoryName{} has again reached a sharper burden than another same-output relay. The current active-primary theorem already showed that the current observer-relevant pico-section minimizer-section core is not primitive: it is the canonical effective output of a deeper current pico-section package \(\PicoSectionPackage\). The accompanying renewed routing note then fixed the consequence of that theorem. The next honest benchmark-facing manuscript must either derive or justify \(\PicoSectionPackage\) itself from still more primitive substrate structure, or prove a genuinely smaller theorem-level collapse strictly inside that package.

The present note takes the direct derivation route. Its claim is not that final substrate microphysics has been identified, and it is not that the benchmark exact-relativistic package has changed. Its claim is narrower and benchmark facing. The current pico-section package is shown to arise canonically from a more primitive femto-section layer whose projected exact-shell transport already contains the benchmark-relevant structure of \(\PicoSectionPackage\). In that precise sense the current pico-section dynamics are no longer the primitive place where the active line has to stop.

Two features matter here. First, the theorem targets the current pico-section package itself rather than re-spending the already settled minimizer-section-core handoff, so it does not count progress merely by reproducing the same current section core from yet another relay. Second, the deeper femto-section package is not merely a renaming of the current pico-section package. It adds an explicit femto-section state space, a femto-section body, a femto-to-pico-section projection, and a deeper fiberwise minimizing law whose effective energy reproduces the recorded pico-section functional on the benchmark-relevant body.

\paragraph{Framework and claim boundary.}
\TheoryProgram{} denotes the broader \Aether{} / \Aflow{} framework built on the claims that the \Aether{} is the underlying four-dimensional substrate of reality, the \Aflow{} is its intrinsic ordered motion, observed three-dimensional space is the local experiential slice of that deeper substrate, S-time is the experienced order of change arising from matter, light, and the \Aflow{}, and observed expansion is the three-dimensional appearance of deeper four-dimensional ordered motion. Within that framework, \TheoryName{} denotes the exact-closure theory in which the effective gravitational dynamics are adopted to be exactly Einsteinian with universal matter coupling. In the active sequence, \emph{adoption} means use of that established relativistic dynamics without claiming substrate derivation, while \emph{derivation} is reserved for a first-principles recovery from explicit substrate variables.

The present note stays strictly inside that boundary. It does not alter the benchmark exact-closure package. It does not introduce a second operative metric or a new low-energy matter law. It only proves that the current nano-section-generating substrate pico-section dynamics can itself be generated from a more primitive femto-section-level substrate dynamics.

\section{Fixed Inputs}

\begin{definition}[Recorded primitive continuation data]
\label{def:fixed}
Fix the following continuation data from the active primitive-reservoir observer-reduction line.
\begin{enumerate}
    \item For each sector label \(\sigma\in\{\Car,\Cur,\Hom\}\), a primitive forcing shell
    \[
    \ForSpace{\sigma}.
    \]
    \item The product shell
    \[
    \PrimShell
    \equiv
    \ForSpace{\Car}\times \ForSpace{\Cur}\times \ForSpace{\Hom}.
    \]
    \item The recorded metric and ordinary-matter outputs
    \[
    \MetricOut:\PrimShell\to\{\text{observer metrics}\},
    \qquad
    \MatterOut:\PrimShell\times\MatterSpace\to\{\text{observer stress tensors}\},
    \]
    satisfying
    \begin{equation}
    \MetricOut(\mathbf w)=\CompletedMetric,
    \qquad
    \MatterOut(\mathbf w,\psi)
    =
    T_{\mu\nu}^{\mathrm{matter}}[\CompletedMetric,\psi]
    \label{eq:fixedoutputs}
    \end{equation}
    for every \(\mathbf w\in\PrimShell\) and every matter configuration \(\psi\in\MatterSpace\).
\end{enumerate}
\end{definition}

\begin{definition}[Recorded pico-section downstream chain]
\label{def:downstream}
Fix \(\sigma\in\{\Car,\Cur,\Hom\}\). The active line already fixes the following downstream data and consequences:
\begin{enumerate}
    \item the recorded nano-section target data
    \[
    \NanoSectionSpace{\sigma},\qquad
    \NanoSectionBody{\sigma},\qquad
    \NanoSectionFunctional{\sigma},\qquad
    \NanoSectionShell{\sigma};
    \]
    \item the current observer-relevant pico-section minimizer-section core \(\PicoMinSectionCore\), the current nano-section package \(\NanoSectionPackage\), the current micro-section package \(\MicroSectionPackage\), the current sub-section package \(\SubSectionPackage\), and the observer-relevant sub-pre-core exact-shell transport-section core \(\SubPreSectionCore\);
    \item the previously recorded fact that, once the current pico-section package is available, the current minimizer-section core, the current nano-section package, the current micro-section package, the current sub-section package, and the downstream shell-transport consequences used on the active line follow unchanged.
\end{enumerate}
\end{definition}

\begin{definition}[Recorded nano-section-generating substrate pico-section dynamics]
\label{def:picosec}
Fix \(\sigma\in\{\Car,\Cur,\Hom\}\). The recorded current pico-section dynamics package \(\PicoSectionPackage\) for sector \(\sigma\) consists of the following data.
\begin{enumerate}
    \item A finite-dimensional Euclidean pico-section state space \(\PicoSectionSpace{\sigma}\), a compact convex pico-section body
    \[
    \PicoSectionBody{\sigma}\subset\PicoSectionSpace{\sigma}
    \]
    with nonempty interior, and an affine surjection
    \[
    \PicoSectionProj{\sigma}:\PicoSectionSpace{\sigma}\to\NanoSectionSpace{\sigma}
    \]
    satisfying
    \[
    \PicoSectionProj{\sigma}\bigl(\PicoSectionBody{\sigma}\bigr)=\NanoSectionBody{\sigma}.
    \]
    For each \(p\in\NanoSectionBody{\sigma}\), write
    \[
    \PicoSectionFiber{\sigma}{p}
    \equiv
    \bigl(\PicoSectionProj{\sigma}\bigr)^{-1}(p)\cap\PicoSectionBody{\sigma}.
    \]
    Each such fiber is required to be nonempty, compact, and convex.
    \item A nonnegative smooth pico-section functional
    \[
    \PicoSectionFunctional{\sigma}:\PicoSectionSpace{\sigma}\to[0,\infty)
    \]
    such that, for every \(p\in\NanoSectionBody{\sigma}\), the restriction of \(\PicoSectionFunctional{\sigma}\) to \(\PicoSectionFiber{\sigma}{p}\) is strictly convex and attains a unique minimizer. The resulting minimizer depends smoothly on \(p\); write it as
    \[
    \PicoSectionNanoMin{\sigma}:\NanoSectionBody{\sigma}\to\PicoSectionBody{\sigma}.
    \]
    These data satisfy
    \[
    \PicoSectionProj{\sigma}\circ\PicoSectionNanoMin{\sigma}
    =
    \mathrm{id}_{\NanoSectionBody{\sigma}},
    \]
    and the effective minimum value recovers the recorded nano-section functional:
    \[
    \NanoSectionFunctional{\sigma}(p)
    =
    \PicoSectionFunctional{\sigma}\bigl(\PicoSectionNanoMin{\sigma}(p)\bigr)
    =
    \min_{r\in\PicoSectionFiber{\sigma}{p}}
    \PicoSectionFunctional{\sigma}(r)
    \qquad
    \text{for every }
    p\in\NanoSectionBody{\sigma}.
    \]
    \item A closed embedded exact pico-section shell
    \[
    \PicoSectionShell{\sigma}\subset\PicoSectionBody{\sigma}
    \]
    satisfying
    \[
    \PicoSectionShell{\sigma}
    =
    \left\{
    r\in\PicoSectionBody{\sigma}:
    \begin{array}{l}
    \PicoSectionProj{\sigma}(r)\in\NanoSectionShell{\sigma},\\[0.2em]
    \PicoSectionFunctional{\sigma}(r)
    =
    \displaystyle\min_{s\in\PicoSectionFiber{\sigma}{\PicoSectionProj{\sigma}(r)}}
    \PicoSectionFunctional{\sigma}(s)
    \end{array}
    \right\},
    \]
    and
    \[
    \PicoSectionProj{\sigma}\big|_{\PicoSectionShell{\sigma}}
    :
    \PicoSectionShell{\sigma}
    \xrightarrow{\cong}
    \NanoSectionShell{\sigma}
    \]
    is a diffeomorphism.
    \item Benchmark faithfulness, meaning preservation of the benchmark outputs \eqref{eq:fixedoutputs}, no second operative metric, no enlarged low-energy matter law, and presentation as a deeper substrate-side source for the current nano-section package rather than as a replacement benchmark theory.
\end{enumerate}
\end{definition}

\begin{remark}
Definitions \ref{def:fixed}--\ref{def:picosec} are fixed inputs. The one operative metric, the ordinary matter route, the recorded pico-section body \(\PicoSectionBody{\sigma}\), the pico-section functional \(\PicoSectionFunctional{\sigma}\), the exact pico-section shell \(\PicoSectionShell{\sigma}\), and the already recorded downstream consequences are not reopened here. The present note asks only whether the current pico-section package can itself be generated from still more primitive substrate structure.
\end{remark}

\section{Pico-Section-Generating Substrate Femto-Section Dynamics}

\begin{definition}[Pico-section-generating substrate femto-section dynamics]
\label{def:femtosec}
Fix \(\sigma\in\{\Car,\Cur,\Hom\}\). A \emph{pico-section-generating substrate femto-section dynamics package} \(\FemtoSectionPackage\) for sector \(\sigma\) consists of the following data.
\begin{enumerate}
    \item A finite-dimensional Euclidean femto-section state space \(\FemtoSectionSpace{\sigma}\), a compact convex femto-section body
    \[
    \FemtoSectionBody{\sigma}\subset\FemtoSectionSpace{\sigma}
    \]
    with nonempty interior, and an affine surjection
    \[
    \FemtoSectionProj{\sigma}:\FemtoSectionSpace{\sigma}\to\PicoSectionSpace{\sigma}
    \]
    satisfying
    \[
    \FemtoSectionProj{\sigma}\bigl(\FemtoSectionBody{\sigma}\bigr)=\PicoSectionBody{\sigma}.
    \]
    For each \(r\in\PicoSectionBody{\sigma}\), write
    \[
    \FemtoSectionFiber{\sigma}{r}
    \equiv
    \bigl(\FemtoSectionProj{\sigma}\bigr)^{-1}(r)\cap\FemtoSectionBody{\sigma}.
    \]
    Each such fiber is required to be nonempty, compact, and convex.
    \item A nonnegative smooth femto-section functional
    \[
    \FemtoSectionFunctional{\sigma}:\FemtoSectionSpace{\sigma}\to[0,\infty)
    \]
    such that, for every \(r\in\PicoSectionBody{\sigma}\), the restriction of \(\FemtoSectionFunctional{\sigma}\) to \(\FemtoSectionFiber{\sigma}{r}\) is strictly convex and attains a unique minimizer. The resulting minimizer depends smoothly on \(r\); write it as
    \[
    \FemtoSectionPicoMin{\sigma}:\PicoSectionBody{\sigma}\to\FemtoSectionBody{\sigma}.
    \]
    These data satisfy
    \[
    \FemtoSectionProj{\sigma}\circ\FemtoSectionPicoMin{\sigma}
    =
    \mathrm{id}_{\PicoSectionBody{\sigma}},
    \]
    and the effective minimum value recovers the recorded pico-section functional:
    \[
    \PicoSectionFunctional{\sigma}(r)
    =
    \FemtoSectionFunctional{\sigma}\bigl(\FemtoSectionPicoMin{\sigma}(r)\bigr)
    =
    \min_{t\in\FemtoSectionFiber{\sigma}{r}}
    \FemtoSectionFunctional{\sigma}(t)
    \qquad
    \text{for every }
    r\in\PicoSectionBody{\sigma}.
    \]
    \item A closed embedded exact femto-section shell
    \[
    \FemtoSectionShell{\sigma}\subset\FemtoSectionBody{\sigma}
    \]
    satisfying
    \[
    \FemtoSectionShell{\sigma}
    =
    \left\{
    t\in\FemtoSectionBody{\sigma}:
    \begin{array}{l}
    \FemtoSectionProj{\sigma}(t)\in\PicoSectionShell{\sigma},\\[0.2em]
    \FemtoSectionFunctional{\sigma}(t)
    =
    \displaystyle\min_{u\in\FemtoSectionFiber{\sigma}{\FemtoSectionProj{\sigma}(t)}}
    \FemtoSectionFunctional{\sigma}(u)
    \end{array}
    \right\},
    \]
    and
    \[
    \FemtoSectionProj{\sigma}\big|_{\FemtoSectionShell{\sigma}}
    :
    \FemtoSectionShell{\sigma}
    \xrightarrow{\cong}
    \PicoSectionShell{\sigma}
    \]
    is a diffeomorphism.
    \item Benchmark faithfulness, meaning preservation of the benchmark outputs \eqref{eq:fixedoutputs}, no second operative metric, no enlarged low-energy matter law, and presentation as a deeper substrate-side source for the current pico-section package rather than as a replacement benchmark theory.
\end{enumerate}
\end{definition}

\begin{remark}
Definition \ref{def:femtosec} is intentionally written over the already fixed pico-section body, pico-section functional, and exact pico-section shell. It does not reopen the already-compressed minimizer-image-core or minimizer-section-core presentations as the live burden. Its positive role is narrower: the current pico-section package is generated as the unique fiberwise effective output of a deeper femto-section selector law.
\end{remark}

\section{Canonical Recovery of the Current Pico-Section Package}

\begin{proposition}[The canonical femto minimizer section is well defined]
\label{prop:section}
For every benchmark-faithful femto-section dynamics package, the map \(\FemtoSectionPicoMin{\sigma}\) of Definition \ref{def:femtosec} is a smooth section of \(\FemtoSectionProj{\sigma}\), so that
\[
\FemtoSectionProj{\sigma}\circ\FemtoSectionPicoMin{\sigma}
=
\mathrm{id}_{\PicoSectionBody{\sigma}}.
\]
\end{proposition}

\begin{proof}
Definition \ref{def:femtosec}(1) makes each fiber \(\FemtoSectionFiber{\sigma}{r}\) nonempty, compact, and convex. Definition \ref{def:femtosec}(2) states that \(\FemtoSectionFunctional{\sigma}\) has a unique minimizer on each such fiber and that the minimizer depends smoothly on \(r\). Definition \ref{def:femtosec}(2) therefore gives a well-defined smooth map
\[
\FemtoSectionPicoMin{\sigma}:\PicoSectionBody{\sigma}\to\FemtoSectionBody{\sigma}.
\]
Because every minimizer lies in the fiber over the point at which it is selected, we have
\[
\FemtoSectionProj{\sigma}\bigl(\FemtoSectionPicoMin{\sigma}(r)\bigr)=r
\]
for every \(r\in\PicoSectionBody{\sigma}\), which is the displayed identity.
\end{proof}

\begin{proposition}[The exact pico-section shell is the projected exact femto-section shell]
\label{prop:shell}
For every benchmark-faithful femto-section dynamics package,
\[
\FemtoSectionProj{\sigma}\bigl(\FemtoSectionShell{\sigma}\bigr)
=
\PicoSectionShell{\sigma},
\]
and the shell restriction
\[
\FemtoSectionProj{\sigma}\big|_{\FemtoSectionShell{\sigma}}
:
\FemtoSectionShell{\sigma}\xrightarrow{\cong}\PicoSectionShell{\sigma}
\]
is a diffeomorphism with inverse
\[
\FemtoSectionPicoMin{\sigma}\big|_{\PicoSectionShell{\sigma}}
:
\PicoSectionShell{\sigma}\xrightarrow{\cong}\FemtoSectionShell{\sigma}.
\]
\end{proposition}

\begin{proof}
Definition \ref{def:femtosec}(3) states directly that
\[
\FemtoSectionProj{\sigma}\big|_{\FemtoSectionShell{\sigma}}
:
\FemtoSectionShell{\sigma}\xrightarrow{\cong}\PicoSectionShell{\sigma}
\]
is a diffeomorphism, so the image identity follows immediately. To identify the inverse, let \(r\in\PicoSectionShell{\sigma}\). By Proposition \ref{prop:section}, \(\FemtoSectionPicoMin{\sigma}(r)\) lies in the fiber over \(r\). Since \(r\in\PicoSectionShell{\sigma}\), Definition \ref{def:femtosec}(3) shows that this point belongs to \(\FemtoSectionShell{\sigma}\). Hence the shell restriction of \(\FemtoSectionPicoMin{\sigma}\) is indeed the inverse diffeomorphism.
\end{proof}

\begin{proposition}[The induced pico-section package is canonical]
\label{prop:canonical}
Fix a benchmark-faithful femto-section dynamics package. If
\[
s:\PicoSectionBody{\sigma}\to\FemtoSectionBody{\sigma}
\]
is any smooth section satisfying
\[
\FemtoSectionProj{\sigma}\circ s
=
\mathrm{id}_{\PicoSectionBody{\sigma}}
\]
and
\[
\FemtoSectionFunctional{\sigma}\bigl(s(r)\bigr)
=
\min_{t\in\FemtoSectionFiber{\sigma}{r}}
\FemtoSectionFunctional{\sigma}(t)
\qquad
\text{for every }
r\in\PicoSectionBody{\sigma},
\]
then \(s=\FemtoSectionPicoMin{\sigma}\). In particular, the induced current pico-section package carries no residual femto-fiber minimizing-choice freedom once the deeper femto-section dynamics are fixed.
\end{proposition}

\begin{proof}
For each \(r\in\PicoSectionBody{\sigma}\), the point \(s(r)\) belongs to \(\FemtoSectionFiber{\sigma}{r}\) and minimizes \(\FemtoSectionFunctional{\sigma}\) there. By Definition \ref{def:femtosec}(2), that minimizer is unique. Therefore
\[
s(r)=\FemtoSectionPicoMin{\sigma}(r)
\]
for every \(r\), hence \(s=\FemtoSectionPicoMin{\sigma}\).
\end{proof}

\begin{proposition}[All current downstream uses factor through the induced pico-section package]
\label{prop:downstream}
For every benchmark-faithful femto-section dynamics package, all current benchmark-facing downstream uses of the current pico-section package \(\PicoSectionPackage\) on the active line factor through the induced pico-section data together with the fixed femto-section package. In particular, the already recorded recovery of the current minimizer-section core \(\PicoMinSectionCore\), the current nano-section package \(\NanoSectionPackage\), the current micro-section package \(\MicroSectionPackage\), the current sub-section package \(\SubSectionPackage\), and the previously recorded shell-transport consequences does not require any additional primitive pico-section datum once the induced pico-section package is fixed.
\end{proposition}

\begin{proof}
By Definition \ref{def:femtosec}(1)--(3), the deeper femto-section package already supplies the recorded pico-section body, pico-section functional, and exact pico-section shell. Proposition \ref{prop:section} gives the canonical femto minimizer section over that body, and Proposition \ref{prop:shell} gives the exact shell projection and inverse lift. Definition \ref{def:downstream}(3) then records that, once the current pico-section package is available, the current minimizer-section core, the current nano-section package, the current micro-section package, the current sub-section package, and the downstream shell-transport consequences used on the active line follow unchanged. Therefore all current benchmark-facing downstream uses at pico-section scope already factor through the induced pico-section package together with the fixed femto-section data, and no additional primitive pico-section datum is required.
\end{proof}

\begin{proposition}[The benchmark package remains fixed]
\label{prop:benchmark}
Every benchmark-faithful pico-section-generating substrate femto-section dynamics package preserves the same benchmark one-metric package \eqref{eq:fixedoutputs}. In particular, the operative metric remains exactly \(\CompletedMetric\), the ordinary matter route is unchanged, there is no second operative metric, and the low-energy matter law is not enlarged.
\end{proposition}

\begin{proof}
This is part of benchmark faithfulness in Definition \ref{def:femtosec}(4). The present note only derives the current pico-section package from deeper substrate selector data. It does not alter the benchmark exact-closure package.
\end{proof}

\section{Main Theorem}

\begin{theorem}[Nano-section-generating substrate pico-section dynamics from pico-section-generating substrate femto-section dynamics]
\label{thm:main}
Fix the primitive continuation data of Definition \ref{def:fixed}, the recorded pico-section downstream chain of Definition \ref{def:downstream}, the recorded current pico-section package of Definition \ref{def:picosec}, and a benchmark-faithful pico-section-generating substrate femto-section dynamics package in the sense of Definition \ref{def:femtosec}. Then, for each sector \(\sigma\in\{\Car,\Cur,\Hom\}\), the current nano-section-generating substrate pico-section dynamics are no longer an unexplained primitive stopping point on the active line: they are recovered canonically from the deeper femto-section dynamics. More precisely:
\begin{enumerate}
    \item the deeper femto-section dynamics generate the current pico-section body, the canonical femto minimizer lift over that body, the effective pico-section functional, and the exact pico-section shell by Definition \ref{def:femtosec}, Proposition \ref{prop:section}, and Proposition \ref{prop:shell};
    \item the induced current pico-section package is canonical and carries no residual femto-fiber minimizing-choice freedom once the deeper package is fixed by Proposition \ref{prop:canonical};
    \item all current benchmark-facing downstream uses of the current pico-section package already factor through the induced datum by Proposition \ref{prop:downstream};
    \item the benchmark boundary remains fixed by Proposition \ref{prop:benchmark};
    \item consequently the remaining honest burden below the renewed pico-section-package frontier is no longer the current nano-section-generating substrate pico-section dynamics \(\PicoSectionPackage\) as such. What remains is to derive or justify the deeper pico-section-generating substrate femto-section dynamics \(\FemtoSectionPackage\) from still more primitive substrate structure, or else prove a still sharper theorem-level collapse strictly inside that deeper package.
\end{enumerate}
\end{theorem}

\begin{proof}
Item (1) is the content of Definition \ref{def:femtosec} together with Propositions \ref{prop:section} and \ref{prop:shell}. Item (2) is Proposition \ref{prop:canonical}. Item (3) is Proposition \ref{prop:downstream}. Item (4) is Proposition \ref{prop:benchmark}. Item (5) is the routing consequence of the previous items together with the active safeguard against counting another same-output deeper-origin relay as new front-line progress when the live benchmark-relevant datum is unchanged.
\end{proof}

\begin{corollary}[What must be done next]
\label{cor:next}
After Theorem \ref{thm:main}, the next honest benchmark-facing manuscript must either:
\begin{enumerate}
    \item derive or justify the pico-section-generating substrate femto-section dynamics \(\FemtoSectionPackage\) from still more primitive substrate structure in a way that changes benchmark-facing status;
    \item identify a genuinely smaller theorem-level observer-relevant datum strictly inside \(\FemtoSectionPackage\) and prove a sharper collapse to it;
    \item do either of the previous two while preserving the same benchmark one-metric package and the same claim boundary.
\end{enumerate}
Another same-output deeper-origin relay that merely preserves the current femto-section package \(\FemtoSectionPackage\) and therefore merely reproduces the same current pico-section package \(\PicoSectionPackage\), the same current minimizer-section core \(\PicoMinSectionCore\), the same current nano-section package \(\NanoSectionPackage\), and the same downstream micro-section and sub-section packages, another re-expansion back to the larger minimizer-image-core or full pico-section presentations as if they were still independently live, or any move that introduces a second operative metric, enlarges the ordinary matter law, or uses stronger promotion language does not satisfy the present burden. If a quantum continuation is pursued, it matters benchmark-facingly only insofar as it derives this same deeper femto-section dynamics, collapses them further, or closes a benchmark derivation gate in a genuinely new way.
\end{corollary}

\begin{proof}
Theorem \ref{thm:main} shows that the current pico-section package is no longer the deepest unexplained datum on the active line. Therefore a subsequent manuscript must improve on the deeper femto-section dynamics rather than merely restate the already generated pico-section package.
\end{proof}

\section{Conclusion}

The positive result of this note is that the current nano-section-generating substrate pico-section dynamics are no longer primitive on the active derivational line. Once one works with a pico-section-generating substrate femto-section dynamics package over the already fixed pico-section body, the current pico-section functional and exact pico-section shell are forced as the effective fiberwise output of a deeper substrate selection law. All current downstream benchmark-facing uses already factor through that induced pico-section package together with the fixed femto-section package.

This preserves the same benchmark one-metric package and the same claim boundary. The result does not yet provide a completed first-principles derivation of GR from substrate microphysics, and it does not license stronger promotion language. Its narrower benchmark-facing gain is that the remaining honest burden now sits one level lower again: not on the current pico-section package itself, but on the deeper femto-section dynamics that canonically generate it, or on a still smaller collapse strictly inside that deeper package.

\end{document}
